\scanpage{7}
On se propose de décrire. \inserted{\deleted{\unsure{canoniquement}} (à l'\unsure{équ.} près...)} $\mathcal M$ \inserted{(les structures gradation, augmentation)} avec la famille des $T_\ell$ ($\ell\in\operatorname{Specmax}(\mathbf Z)$), \deleted{le foncteur $\unsure{T_\infty}$ et} et le foncteur $F_0$. La connaissance de $\mathcal M$, gradation, augmentation, $F_0$ est connue quand on connaît la classe de $\mathcal H$ comme $H$-gerbe, i.e. un élément
\[
\xi_0\in H^2(Q,H)\simeq H^2(Q,H^0)\times H^2(Q,\mu_2)
\]
N.B. la gradation correspond à un hom
\[
\Gamma\xrightarrow{i^*}\mathbf Z
\]
qui est donné par
\[
i^*(\lambda)=\nu\quad\Longleftrightarrow\quad |\lambda|=p^{\nu/2}).
\]
On voit :

Si $\ell\ne p$, le foncteur $F_\ell$ prolonge (à isom près) le foncteur $F_0$, donc \deleted{\unsure{est}} définit un splitting de la gerbe $\mathcal H$ \deleted{\unsure{en}} en $\ell$. \uline{Donc} $\mathcal H$ \uline{est splittée canoniquement en les} $\ell\ne p$ [i.e. sur les $Q_\ell$]

donc la connaissance de \deleted{\unsure{la gerbe et}} $\mathcal M,\mathrm{grad},\mathrm{aug},F_0,F_\ell\ (\ell\ne p)$ équivaut à la donnée de
\[
\xi\in H^2(\unsure{S|S_0}\bmod T,H)
\]
où $T=\coprod_{\ell\ne p}\operatorname{Spec}(Q_\ell)$. Retenons pour l'instant que $\xi\in H^2(Q,H)$ \uline{splitte en les places finies} $\ell\ne p$.

Si $\ell=p$, les foncteurs fibres $T_p\!\mid\!\operatorname{Niv}^0(\mathcal M)^0$ et $\unsure{F_p|F_0}$ ne sont pas isomorphiques. Ils diffèrent par un torseur sous $\unsure{(G/H)_{Q_p}|(G/H)_{\overline Q_p}}$ (extension de $G/G^0\simeq\widehat{\mathbf Z}$ par $G^0/H\simeq G_m$,
\mspage{6}
